\frfilename{mt11.tex}
\versiondate{4.1.04}
\copyrightdate{1994}

\def\chaptername{Ruang ukur}
\def\sectionname{Pendahuluan}

\newchapter{11}

Dalam bab ini saya menguraikan konsep dasar `ruang ukur', yaitu suatu
himpunan yang sebagian (biasanya bukan semua) himpunan bagiannya dapat
diberi suatu `ukuran'. Anda dapat menafsirkan ukuran itu sebagai luas,
massa, volume, kapasitas termal, atau hampir apa saja yang Anda harapkan
bersifat aditif -- maksud saya, ukuran gabungan dua himpunan yang saling
lepas harus sama dengan jumlah ukuran keduanya. Definisi sebenarnya (dalam
112A) tidaklah langsung terlihat dan pada hakikatnya bergantung pada
sejumlah ciri teknis yang membuat sebuah bagian persiapan (\S111) layak
disajikan. Selain itu, sekalipun definisinya sudah dipahami dengan baik,
contoh ukuran yang paling awal dan paling penting, yakni ukuran Lebesgue
pada ruang Euklides, tetap tidak mudah ditangkap. Karena itu, sebelum
beralih ke perincian argumen yang diperlukan untuk ukuran Lebesgue dalam
\S\S114-115, saya mencurahkan satu bagian (\S113) untuk suatu metode
membangun ukuran. Jadi, struktur bab ini terdiri atas tiga bagian teori
umum, disusul dua bagian (yang sangat mirip) mengenai contoh-contoh
khusus. Perlu saya katakan bahwa teori umumnya pada dasarnya lebih mudah;
namun teori itu memang mengandalkan kelancaran dalam melakukan manipulasi
tertentu terhadap keluarga himpunan yang mungkin masih baru bagi Anda.

Pada suatu saat saya perlu menjelaskan penataan materi ini, dan mungkin
akan membantu jika saya melakukannya sebelum Anda mulai mempelajari bab
ini. Salah satu dari banyak pertanyaan mendasar yang harus diputuskan oleh
setiap penulis mengenai pokok ini ialah apakah akan memulai dengan teori
ukur `umum' atau dengan ukuran dan integrasi Lebesgue. Masalahnya, ukuran
Lebesgue bukan sekadar contoh ruang ukur yang paling penting. Ukuran itu
begitu dekat dengan inti pokok bahasan sehingga sebagian besar gagasan
dalam teori elementer dapat diwujudkan sepenuhnya dalam teorema-teorema
tentang ukuran Lebesgue. Jika menengok ke Jilid 2, saya mendapati bahwa,
selain Bab 21 -- yang secara khusus dimaksudkan untuk memperluas gagasan
Anda mengenai apa saja yang dapat menjadi ruang ukur -- hanya Bab 27 dan
paruh kedua Bab 25 yang benar-benar memerlukan teori umum agar dapat
dipahami, sedangkan Bab 22, 26, dan 28 secara khusus membahas ukuran
Lebesgue. Jilid 3 lain lagi persoalannya, tetapi bahkan di sana lebih dari
separuh kandungan matematisnya dapat dinyatakan dengan ukuran Lebesgue.
Jika Anda berpandangan, sebagaimana tentu saya lakukan ketika pandangan
itu sesuai dengan argumen saya, bahwa tugas matematika murni ialah
mengungkapkan dan memperluas kemampuan logis pikiran manusia, sedangkan
teorema-teorema yang kita pelajari hanyalah wahana bagi gagasan-gagasan itu, maka
Anda dapat dengan tepat menunjukkan bahwa semua hal yang sungguh penting
dalam jilid ini dapat dilakukan tanpa repot-repot merumuskan teori umum
ruang ukur abstrak; dan bahwa dengan mempelajari contoh yang relatif
konkret berupa ukuran Lebesgue pada ruang Euklides berdimensi $r$, Anda
dapat menghindari berbagai gangguan yang tidak relevan.

Jika, sebagai pengajar, Anda benar-benar yakin bahwa tidak seorang pun
dari peserta didik Anda ingin melampaui teori elementer, pandangan ini
patut dipertimbangkan. Namun, saya percaya bahwa pandangan itu tidak dapat
dipertahankan jika Anda ingin mempersiapkan seorang pun di antara peserta
didik Anda untuk gagasan-gagasan yang lebih maju. Kesulitannya ialah bahwa, betapapun baik
niatnya, orang yang telah mempelajari teori lengkap ukuran Lebesgue lalu
beralih ke teori ruang ukur abstrak cenderung melaluinya terlalu cepat,
dan pada akhirnya tidak lagi yakin persis fakta-fakta mana—yang jumlahnya
sembilan puluh persen dari semua fakta yang diketahuinya—yang berlaku secara
umum. Saya percaya akan lebih aman jika
sifat-sifat khusus ukuran Lebesgue sejak awal diberi label yang jelas
sebagai sifat khusus.

Tentu saja, kekeliruan yang terus menghantui pengajar matematika pada
tingkat ini ialah mengajar kelas berisi dua puluh orang dengan cara yang
mungkin hanya cocok bagi dua orang di antaranya. Namun, dalam hal ini
penilaian saya sendiri ialah bahwa sangat sedikit mahasiswa yang memang
siap mengikuti mata kuliah ini akan mengalami kesulitan dengan tingkat
abstraksi tambahan dalam `Misalkan $(X,\Sigma,\mu)$ suatu ruang ukur,
$\ldots$'. Saya memang mengandaikan pengetahuan tentang aljabar linear
elementer, dan setidaknya tata bahasa ruang ukur sembarang tidak lebih
buruk daripada tata bahasa ruang linear sembarang. Lagi pula, teori
Lebesgue sudah melibatkan pernyataan berbentuk `jika $E$ merupakan
himpunan terukur Lebesgue, $\ldots$', dan menurut pengalaman saya,
mahasiswa yang mampu menangani kuantifikasi atas himpunan bagian bilangan
real tidak gentar menghadapi kuantifikasi atas himpunan yang
unsur-unsurnya juga himpunan (yang bagaimanapun diperlukan dalam setiap
uraian elementer tentang aljabar-$\sigma$ himpunan Borel). Jadi saya
percaya bahwa, setidaknya di sini, keumuman tambahan dari pendekatan
`profesional' bukanlah penghalang bagi peminat nonprofesional.

Saya menuliskan semua ini di sini, bukan nanti di dalam bab, karena saya
memang ingin memberi Anda pilihan. Jika Anda memilih mempelajari teori
Lebesgue terlebih dahulu dan menunda teori umum, lakukanlah sebagai
berikut. Anda perlu membaca

\qquad paragraf 114A-114C %114A 114B 114C

\qquad 114D, bersama 113A-113B dan 112Ba, 112Bc

\qquad 114E, bersama 113C-113D, 111A, 112A, 112Bb

\qquad 114F

\qquad 114G, bersama 111G dan 111C-111F, %111C 111D 111E 111F

\noindent lalu lanjutkan dengan Bab 12. Pada suatu saat, tentu saja, Anda
perlu melihat latihan-latihan untuk \S\S112-113; tetapi, seperti dalam Bab
12-13, lakukanlah dengan menerjemahkan `Misalkan $(X,\Sigma,\mu)$ suatu
ruang ukur' menjadi `Misalkan $\mu$ adalah ukuran Lebesgue pada $\Bbb R$,
dan $\Sigma$ aljabar-$\sigma$ himpunan terukur Lebesgue'. Demikian pula,
ketika melihat 111X-111Y, ambillah $\Sigma$ sebagai {\it entah}
aljabar-$\sigma$ himpunan terukur Lebesgue {\it atau} aljabar-$\sigma$
himpunan bagian Borel dari $\Bbb R$.

\discrpage
